\documentclass[addpoints]{exam}
% \documentclass[addpoints,12pt]{exam}

\usepackage[slovene]{babel}
\usepackage{amsfonts}
\usepackage{amssymb}
\usepackage{amsmath}
\usepackage{float}
\usepackage{graphicx}
\usepackage{enumitem}
\usepackage{color}
\usepackage{eurosym}


%Komentar naslednje vrstice glede na verzijo testa -- rešitve ali prostor za odgovore:
% \printanswers

% \linespread{2}


\pointsinrightmargin
\marginbonuspointname{}


\renewcommand{\solutiontitle}{\noindent\textbf{Rešitve in točkovanje:}\par\noindent}

\colorgrids
\definecolor{GridColor}{gray}{0.7}
\setlength{\gridsize}{7mm}

\extrawidth{5mm}
\setlength{\rightpointsmargin}{1.5cm}

\begin{document}

\runningfooter{}{}{}

\begin{center}
    \fbox{\fbox{\parbox{15cm}{\centering
    Prvo pisno ocenjevanje znanja \\ 1. letnik -- test C \\ 16. oktober 2025 \\ (Osnove logike in teorije množic, Naravna in cela števila)}}}
\end{center}

\vspace{0.1in}
\makebox[\textwidth]{Ime in priimek, oddelek:\enspace\hrulefill}


\begin{center}
    \chqword{Naloga}
    \chpword{Možne točke}
    \chbpword{Dodatne točke}
    \chsword{Dosežene točke}
    \chtword{Skupaj}  
    % \combinedpointtable[h][questions]
    \combinedgradetable[h][questions]

\end{center}

\vspace{0.1in}


\begin{questions}

    % 1. vprašanje
    
    \question[4]
    Določite pravilnost izjave $(\lnot A \lor  B)\Rightarrow (A\Leftrightarrow \lnot B)$ glede na izbor izjav $A$ in $B$.

    \begin{solution}[11cm]
        \begin{table}[H]
            \begin{tabular}{ccccccc}
             $A$ & $B$ & $\lnot A$ & $\lnot B$ & $(\lnot A\Rightarrow B)$ & $(A\Leftrightarrow \lnot B)$ & $(\lnot A\Rightarrow B)\land(A\Leftrightarrow \lnot B)$\\ \hline
             P & P & N & N & P & N & N \\ 
             P & N & N & P & P & P & P \\
             N & P & P & N & P & P & P \\
             N & N & P & P & N & N & N
            \end{tabular}
        \end{table}
        Za vsako pravilno izpolnjeno vrstico, ki določa pravilnost, po $1$ točka.
    \end{solution}




    % 2. vprašanje
    \question[3]
    Določite resničnostno vrednost izjav:
    \begin{parts}
        \part 
        Če je sodih števil končno mnogo, potem je množica naravnih števil števno neskončna. \\
        \part 
        Število $1$ je praštevilo ali število $5$ je večkratnik števila $1$.\\
        \part
        Ni res, da $0$ ni najmanjše celo število, in $0$ je najmanjše naravno število.\\
    \end{parts}

    \begin{solution}
        (a): N (obe izjavi sta nepraavilni) \\
        (b): P (obe izjavi sta pravilni) \\
        (c): P (negirana prva izjava je pravilna, druga izjava je pravilna) \\ \\
        Za vsako pravilno določeno logično vrednost po $1$ točka.
    \end{solution}



    \newpage

    % 3. vprašanje
    \question
    Dani sta množici $\mathcal{A}=\left\{R, A, D, I, O\right\}$ in $\mathcal{B}=\left\{R, I, O\right\}$.
    \begin{parts}
        \part[2]
        Ali velja $\mathcal{B}\subseteq\mathcal{A}$? Utemeljite. 
        \part[1]
        Poiščite najmanjšo možno množico $\mathcal{C}$, da bosta množici $\mathcal{A}$ in $\mathcal{B}$ njeni podmnožici.
        \part[3]
        Zapišite $\mathcal{PB}$. Izračunajte, koliko elementov vsebuje.
        \part[3]
        Zapišite in grafično predstavite kartezični produkt $\mathcal{A}\times\mathcal{B}$.
    \end{parts}

    % \begin{description}
    %     \item[a] $(x-a)^2+(y-b)^2-c^2=0$ 
    %     \item[b] $d^2(x-a)^2+f^2(y-b)^2-(df)^2=0$
    %     \item[c] $S(a,b)$
    %     \item[d] $e^2=d^2-f^2;\ d>f$
    % \end{description}
    % Krožnica: \fillin[a c][2cm]
    % \\ Elipsa:  \fillin[b c d][2cm]

    \begin{solution}[9cm]
        (a): Drži, da $\mathcal{A}\nsubseteq\mathcal{B}$, ker $a\in\mathcal{A}\land a\notin\mathcal{B}$.; $\mathcal{C}=\left\{b,c,h\right\}$ \\
        (b): $\mathcal{PB}=\left\{\emptyset, \{b\}, \{c\}, \{h\}, \{i\}, \{b,c\}, \{b,h\}, \{b,i\}, \{c,h\}, \{c,i\}, \{h,i\}, \{b,c,h\}, \{b,c,i\}, \{b,h,i\}, \{c,h,i\}, \mathcal{B}\right\}$; $m(\mathcal{PB})=2^{m(\mathcal{B})}=2^4=16$. \\
        (c): $\mathcal{B}\times\mathcal{A}=\left\{(a,b), (a,c), (a,h), (a,i), (b,b), (b,c), (b,h), (b,i), (c,b), (c,c), (c,h), (c,i), (h,b), (h,c), (h,h), (h,i)\right\}$ + grafična upodobitev \\ \\
        Pri (a) za vsak odgovor po $1$ točka. Pri (b) za zapis celotne $\mathcal{PB}$ $2$ točki, za zapis zgolj podmnožic $1$ točka; za izračun po formuli za moč potenčne množice $1$ točka. Pri (c) za zapis $\mathcal{B}\times\mathcal{A}$ $1$ točka, za grafično upodobitev $1$ točka.
    \end{solution}


    \newpage
    % 4. vprašanje  
    \question
    Množica $\mathcal{A}$ je prava podmnožica množice $\mathcal{B}$, prav tako je $\mathcal{C}$ prava podmnožica množice $\mathcal{B}$. Množici $\mathcal{A}$ in $\mathcal{C}$ sta disjunktni.
    \begin{parts}
        \part[2]
        Grafično predstavite množico $\mathcal{B}\setminus\left(\mathcal{A}\cup\mathcal{C}\right)$.
        \part[3]
        Izračunajte $\mathcal{A}\setminus\mathcal{C}$, $\mathcal{B}\cap\mathcal{C}$ in $\mathcal{A}\cup\mathcal{B}$.
    \end{parts}

    \begin{solution}[6cm]
        (a): grafična upodobitev \\
        (b): $\mathcal{B}\setminus\mathcal{C}=\emptyset$; $\mathcal{A}\cap\mathcal{C}=\mathcal{A}$; $\mathcal{A}\cup\mathcal{B}=\mathcal{B}$ \\ \\
        Pri (a) za narisan Euler-Vennov diagram $1$ točka. Pri (b) za izračun vsake izmed množic po $1$ točka.
    \end{solution}






    
    % 6. vprašanje
    \question
    Pri uri geografije so se dijaki in učitelj pogovarjali o atlasih.  Učitelj je dijake vprašal, katere atlase imajo doma.
    $20$ dijakov ima \textit{Atlas sveta}, $8$ jih ima \textit{Atlas Jugoslavije} in $17$ \textit{Atlas Slovenije}. Kar $12$ jih ima \textit{Atlas sveta} in \textit{Atlas Slovenije}, $4$ imajo \textit{Atlas Jugoslavije} in \textit{Atlas sveta},
    $3$ pa \textit{Atlas Jugoslavije} in \textit{Atlas Slovenije}. En sam dijaki ima vse tri atlase, dva pa nimata doma nobenega atlasa. 
    \begin{parts}
        \part[1]
        Narišite diagram.
        \part[2]
        Koliko dijakov je bilo pri uri geografije?
        \part[2]
        Koliko dijakov ima doma vsaj dva različna atlasa?
    \end{parts}

    \begin{solution}[8cm]
        (a): Euler-Vennov diagram z vpisanimi številkami \\
        (b): $5+0+3+3+2+11+1+1=26$ dijakov \\
        (c): $3+2+11=16$ dijakov \\ \\
        Pri (a) za pravilno narisan in izpolnjen diagram $1$ točka. Pri (b) in (c) za pravilen izračun $1$ točka, za zapis odgovora $1$ točka.
    \end{solution}




    \newpage

    % 5. vprašanje
    \question
    Naj bo $\mathcal{U}=\mathbb{N}_{21}$. Dane so množice $\mathcal{A}=\left\{x;\ x=2n \land 4\leq n \leq10\right\}$, $\mathcal{B}=\left\{x;\ x=4k-3 \land 3\leq k< 6\right\}$ in $\mathcal{C}=\left\{x;\ x=5m \land 1<m\leq 4\right\}$. 
    \begin{parts}
        \part[3]
        Zapišite dane množice z naštevanjem elementov.
        \part[4]
        Določite naslednje množice: $\mathcal{A}\cap\mathcal{C}$, $\mathcal{C}\setminus\mathcal{A}$, $(\mathcal{C}\cup\mathcal{A})\setminus\mathcal{B}$ in $(\mathcal{A}\cup\mathcal{B})^\complement$.
        \part[1]
        Koliko elementov ima množica $\mathcal{A}\times\mathcal{C}$?
    \end{parts}

    \begin{solution}[13cm]
        $\mathcal{A}=\left\{6,8,10,12,14,16,18\right\}$, $\mathcal{B}=\left\{8,11,14,17,20\right\}$, $\mathcal{C}=\left\{10,15,20\right\}$ \\
        (a): $\mathcal{A}\cap\mathcal{C}=\left\{10\right\}$; $\mathcal{C}\setminus\mathcal{A}=\left\{15,20\right\}$; $(\mathcal{A}\cup\mathcal{C})\setminus\mathcal{B}=\left\{6,10,12,15,16,18\right\}$; $(\mathcal{A}\cup\mathcal{B})^\complement=\left\{1,2,3,4,5,7,9,13,15,19,21,22,23\right\}$ \\
        (b): $m(\mathcal{A}\times\mathcal{C})=m(\mathcal{A})\cdot m(\mathcal{C})=7\cdot 3 = 21$ \\ \\
        Pri (a) za zapis/uporabo pravilnih množic $\mathcal{A}, \mathcal{B}, \mathcal{C}$ $1$ točka; za vsako pravilno določeno množico po $1$ točka. Pri (b) za uporabo formule in izračun moči kartezičnega produkta $1$ točka.
    \end{solution}


    % \newpage
        

    % 7. vprašanje
    \question[3]
    Na smučarskem taboru so oblikovali več različno velikih skupin tečaja smučanja.
    V pet skupin je bilo razporejenih po $12$ tečajnikov, v osem skupin po $9$ tečajnikov, štiri skupine so štele po $7$ tečajnikov, $22$ skupin pa je bilo individualnih.
    V tednu, ko naj bi bil tečaj smučanja izveden, je zbolelo $8$ parov, ki so bili vključeni v tečaj, ostali so se tečaja udeležili.
    Koliko tečajnikov se je udeležilo tečaja?

    \begin{solution}[5cm]
        $5\cdot 26+7\cdot 27+4\cdot 28-320=111$ \\ \\
        Za nastavek enega računa $1$ točka, za pravilen izračun $1$ točka, za zapis odgovora $1$ točka.
    \end{solution}


\newpage
        % dodatno vprašanje

    \makebox[\textwidth]{Ime in priimek, oddelek:\enspace\hrulefill}
    \vspace{0.1in}

    \bonusquestion[3]
    \textit{Dodatna naloga:} Ovrednotite pravilnost izjav: $A: \emptyset\in\mathbb{N}$; $B:\left\{1\right\}\in\mathbb{N}$; $C: \left\{1\right\}\subseteq\left\{\mathbb{N}\right\}$. Odgovore utemeljite.
    \begin{solution}
        $\exists x\in\mathbb{Z}\ni:(3\mid x)\land(2\mid x)$; negacija: $\forall x\in\mathbb{Z}: (3\nmid x)\lor(2\nmid x)$ \\ \\
        Za pravilen zapis izjave s simboli $1$ točka, za pravilno negacijo kvantifikatorja $1$ točka, za pravilno negacijo logične operacije $1$ točka.
    \end{solution}

\end{questions}



\end{document}